\part{Entscheiden für andere}

\chapter{Wenn mehrere entscheiden}

\section{Drei Ärzte, drei Rangfolgen}

\begin{merksatz}
\boxtitel{Die Lage}
Eine Konferenz soll zwischen drei Behandlungen $A$, $B$, $C$ wählen. Die Stimmen verteilen sich so:

\begin{center}
{\small
\renewcommand{\arraystretch}{1.5}
\begin{tabular}{@{}P{2.0cm}P{2.4cm}P{2.4cm}@{}}
\toprule
\textbf{Anteil} & \textbf{Rangfolge} & \\
\midrule
$35\,\%$ & $A \succ B \succ C$ & \\
$33\,\%$ & $B \succ C \succ A$ & \\
$32\,\%$ & $C \succ A \succ B$ & \\
\bottomrule
\end{tabular}}
\end{center}

\schritt{Die Paarvergleiche}
\begin{center}
{\small
\renewcommand{\arraystretch}{1.48}
\begin{tabular}{@{}P{2.4cm}P{2.6cm}P{2.4cm}@{}}
\toprule
\textbf{Vergleich} & \textbf{Ergebnis} & \\
\midrule
$A$ gegen $B$ & $67 : 33$ & $A$ gewinnt\\
$B$ gegen $C$ & $68 : 32$ & $B$ gewinnt\\
$C$ gegen $A$ & $65 : 35$ & \textcolor{akzent}{$C$ gewinnt}\\
\bottomrule
\end{tabular}}
\end{center}

\begin{center}
\begin{minipage}{0.9\textwidth}
\textbf{$A$ schlägt $B$, $B$ schlägt $C$, $C$ schlägt $A$ --- jeweils mit großer Mehrheit.}
\end{minipage}
\end{center}

\begin{center}
\begin{minipage}{0.88\textwidth}
\kursiv{\textbf{Jeder einzelne Arzt hat eine widerspruchsfreie Rangfolge; die Gruppe hat keine.} Wer die Abstimmungsreihenfolge festlegt, bestimmt das Ergebnis: Stimmt man erst über $A$ gegen $B$ ab und den Sieger gegen $C$, gewinnt $C$. Beginnt man mit $B$ gegen $C$, gewinnt $A$. Die Tagesordnung ist dann wichtiger als die Meinungen.}
\end{minipage}
\end{center}
\end{merksatz}

\section{Satz von Arrow}

\begin{merksatz}
\boxtitel{Satz 13.1 --- Arrow, 1951}
Bei mindestens drei Möglichkeiten gibt es kein Verfahren, das aus beliebigen Einzelrangfolgen eine Gruppenrangfolge macht und zugleich erfüllt:

\begin{center}
{\small
\renewcommand{\arraystretch}{1.55}
\begin{tabular}{@{}P{2.8cm}P{5.8cm}@{}}
\toprule
\textbf{Uneingeschränkt} & Jede Zusammenstellung von Einzelrangfolgen ist zugelassen.\\
\textbf{Widerspruchsfrei} & Die Gruppenrangfolge ist vollständig und übertragbar.\\
\textbf{Einstimmigkeit} & Ziehen alle $X$ dem $Y$ vor, so auch die Gruppe.\\
\textbf{Unabhängigkeit} & Die Rangfolge von $X$ und $Y$ hängt nur davon ab, wie die Einzelnen $X$ und $Y$ vergleichen.\\
\midrule
\textbf{Kein Diktator} & Es gibt niemanden, dessen Rangfolge stets durchdringt.\\
\bottomrule
\end{tabular}}
\end{center}

\begin{center}
\begin{minipage}{0.9\textwidth}
\textbf{Die ersten vier Bedingungen erzwingen den Diktator. Es gibt keine Ausnahme.}
\end{minipage}
\end{center}

\begin{center}
\begin{minipage}{0.88\textwidth}
\kursiv{Der Beweis ist nicht lang, aber kunstvoll; er konstruiert aus den Bedingungen schrittweise eine \enquote{entscheidende} Person. Wir übergehen ihn --- die erste Auslassung dieses Bandes; die Fundstelle ist Gaertner\zit{9}. \textbf{Das ist das dritte überraschende Ergebnis des Vorworts}, und es erklärt den Zirkel aus Abschnitt 13.1 nicht als Zufall, sondern als Notwendigkeit.}
\end{minipage}
\end{center}
\end{merksatz}

\begin{merksatz}
\boxtitel{Wie man dennoch entscheidet}
Jede brauchbare Lösung gibt eine der Bedingungen auf.

\begin{center}
{\small
\renewcommand{\arraystretch}{1.55}
\begin{tabular}{@{}P{2.8cm}P{5.8cm}@{}}
\toprule
\textbf{Mehr als Rangfolgen} & Man erhebt Nutzenwerte statt bloßer Reihungen und addiert sie. Gibt die Unabhängigkeit auf.\\
\textbf{Eingeschränkte Vorlieben} & Lassen sich alle Möglichkeiten auf einer Linie ordnen und hat jeder ein Optimum darauf, so entscheidet die Mitte widerspruchsfrei.\\
\textbf{Ein Entscheider} & In der Klinik trägt am Ende einer die Verantwortung --- der \enquote{Diktator} ist hier kein Schimpfwort, sondern eine Zuständigkeit.\\
\bottomrule
\end{tabular}}
\end{center}

\begin{center}
\begin{minipage}{0.88\textwidth}
\kursiv{\textbf{Die zweite Zeile ist für die Medizin die wichtigste}: Behandlungen lassen sich oft von \enquote{zurückhaltend} bis \enquote{aggressiv} ordnen, und jeder Beteiligte hat darauf einen bevorzugten Punkt. Dann gibt es einen Gewinner, der gegen jeden anderen eine Mehrheit findet --- den Standpunkt der Mitte. Konferenzen arbeiten faktisch so, ohne es zu wissen.}
\end{minipage}
\end{center}
\end{merksatz}

\section{Der Bruchfall}

\begin{warnung}
\boxtitelr{Was Nutzenaddition voraussetzt}
\begin{center}
{\small
\renewcommand{\arraystretch}{1.52}
\begin{tabular}{@{}P{2.8cm}P{5.8cm}@{}}
\toprule
\textbf{Vergleichbarkeit} & Nach Satz 2.1 ist jeder Nutzen nur bis auf Nullpunkt und Maßstab bestimmt --- ein Vergleich zwischen Personen setzt also eine zusätzliche Festlegung voraus.\\
\textbf{Übertreibung} & Wer seinen Nutzen selbst angibt, hat Anlass, ihn zu übertreiben --- eine Frage aus Band 21, Kapitel~13.\\
\textbf{Minderheiten} & Eine Summe kann eine kleine Gruppe beliebig schlecht stellen, wenn die Mehrheit genug gewinnt.\\
\bottomrule
\end{tabular}}
\end{center}

\begin{center}
\begin{minipage}{0.88\textwidth}
\kursiv{\textbf{Die erste Zeile ist grundsätzlich}: Dass Frau B.\ einen Nutzenwert von $0{,}65$ angibt und Herr K.\ einen von $0{,}40$, sagt nichts darüber, wer mehr leidet --- die Skalen sind je eigene. Die Gesundheitsökonomie legt deshalb eine gemeinsame Skala fest: null gleich Tod, eins gleich volle Gesundheit für alle. \textbf{Das ist eine Setzung, keine Messung}, und Kapitel~14 zeigt, was sie kostet.}
\end{minipage}
\end{center}
\end{warnung}

\section{Übungsaufgaben}

\begin{aufgabe}
\boxtitel{Aufgabe 13.1 --- Die Tagesordnung entscheidet}
Zeigen Sie an den Zahlen aus Abschnitt 13.1, dass jede der drei Behandlungen gewinnen kann.

\loesung
\schritt{Schritt 1 --- Die drei Reihenfolgen}
\begin{center}
{\small
\renewcommand{\arraystretch}{1.48}
\begin{tabular}{@{}P{3.0cm}P{2.4cm}P{2.2cm}@{}}
\toprule
\textbf{zuerst} & \textbf{dann gegen} & \textbf{Sieger}\\
\midrule
$A$ gegen $B$ ($A$) & $C$ & $\mathbf{C}$\\
$B$ gegen $C$ ($B$) & $A$ & $\mathbf{A}$\\
$C$ gegen $A$ ($C$) & $B$ & $\mathbf{B}$\\
\bottomrule
\end{tabular}}
\end{center}

\schritt{Schritt 2 --- Die Folge}
\begin{center}
\begin{minipage}{0.9\textwidth}
\textbf{Wer die Tagesordnung setzt, wählt den Sieger --- ohne eine einzige Stimme zu ändern.}
\end{minipage}
\end{center}

\schritt{Schritt 3 --- Was hilft}
\begin{center}
{\small
\renewcommand{\arraystretch}{1.5}
\begin{tabular}{@{}P{2.6cm}P{6.0cm}@{}}
\toprule
\textbf{Alle Paare zeigen} & Der Zirkel wird sichtbar statt versteckt.\\
\textbf{Punkte vergeben} & Nach Borda: $2$, $1$, $0$ je Rang. Hier: $A$ $102$, $B$ $101$, $C$ $97$ --- eindeutig, aber abhängig davon, welche Möglichkeiten zur Wahl stehen.\\
\bottomrule
\end{tabular}}
\end{center}

\begin{center}
\begin{minipage}{0.88\textwidth}
\kursiv{\textbf{Das Punkteverfahren liefert immer ein Ergebnis und verletzt dafür die Unabhängigkeit}: Streicht man eine chancenlose Möglichkeit, so kann sich die Reihenfolge der übrigen umkehren. Genau dieselbe Schwäche hatte die Minimax-Reue in Kapitel~11. Satz 13.1 sagt, dass man dieser Wahl nicht entgehen kann.}
\end{minipage}
\end{center}
\end{aufgabe}

\begin{aufgabe}
\boxtitel{Aufgabe 13.2 --- Vorlieben auf einer Linie}
Drei Ärzte bevorzugen auf der Achse \enquote{zurückhaltend bis aggressiv} die Punkte $2$, $5$ und $9$. Zur Wahl stehen die Stufen $2$, $5$ und $9$. Gibt es einen Zirkel?

\loesung
\schritt{Schritt 1 --- Die Rangfolgen}
Jeder ordnet nach dem Abstand zu seinem Optimum:
\begin{center}
{\small
\renewcommand{\arraystretch}{1.48}
\begin{tabular}{@{}P{1.8cm}P{4.2cm}@{}}
\toprule
\textbf{Arzt} & \textbf{Rangfolge}\\
\midrule
$2$ & $2 \succ 5 \succ 9$\\
$5$ & $5 \succ 2 \succ 9$\\
$9$ & $9 \succ 5 \succ 2$\\
\bottomrule
\end{tabular}}
\end{center}
\kursiv{Für den mittleren Arzt: Abstand zu $2$ ist $3$, zu $9$ ist $4$ --- also $2$ vor $9$.}

\schritt{Schritt 2 --- Die Paarvergleiche}
\begin{center}
{\small
\renewcommand{\arraystretch}{1.48}
\begin{tabular}{@{}P{2.2cm}P{2.4cm}P{2.2cm}@{}}
\toprule
\textbf{Vergleich} & \textbf{Stimmen} & \textbf{Sieger}\\
\midrule
$5$ gegen $2$ & $2 : 1$ & $\mathbf{5}$\\
$5$ gegen $9$ & $2 : 1$ & $\mathbf{5}$\\
$2$ gegen $9$ & $2 : 1$ & $2$\\
\bottomrule
\end{tabular}}
\end{center}

\begin{center}
\begin{minipage}{0.9\textwidth}
\textbf{Kein Zirkel: $5$ schlägt beide anderen. Die Rangfolge lautet $5 \succ 2 \succ 9$.}
\end{minipage}
\end{center}

\schritt{Schritt 3 --- Warum das immer gilt}
\begin{center}
\begin{minipage}{0.88\textwidth}
\kursiv{\textbf{Der mittlere Standpunkt gewinnt jeden Paarvergleich}: Bei jedem Vergleich stimmt der mittlere Arzt für die ihm nähere Stufe, und mindestens einer der beiden anderen tut dasselbe --- nämlich der, auf dessen Seite sie liegt. Das ist der Satz vom Medianwähler; er stammt von \Person{Duncan Black} aus dem Jahr 1948 und ist der wichtigste Ausweg aus Satz 13.1.}
\end{minipage}
\end{center}

\schritt{Schritt 4 --- Und wann er versagt}
\begin{center}
\begin{minipage}{0.88\textwidth}
\kursiv{Er versagt, sobald jemand zwei getrennte Gipfel hat --- etwa ein Arzt, der entweder abwarten oder sofort operieren will, aber die halbe Maßnahme für das Schlechteste hält. \textbf{Solche Vorlieben sind in der Medizin häufig}, weil halbe Maßnahmen Risiken beider Wege verbinden können. Dann kehrt der Zirkel zurück.}
\end{minipage}
\end{center}
\end{aufgabe}

\begin{aufgabe}
\boxtitel{Aufgabe 13.3 --- Nutzen addieren}
Für dieselben drei Behandlungen geben die Ärzte Nutzenwerte an. Was ändert das?

\begin{center}
{\small
\renewcommand{\arraystretch}{1.48}
\begin{tabular}{@{}P{2.0cm}P{1.8cm}P{1.8cm}P{1.8cm}@{}}
\toprule
\textbf{Arzt} & $A$ & $B$ & $C$\\
\midrule
I ($35$) & $1{,}00$ & $0{,}90$ & $0{,}00$\\
II ($33$) & $0{,}00$ & $1{,}00$ & $0{,}95$\\
III ($32$) & $0{,}40$ & $0{,}00$ & $1{,}00$\\
\bottomrule
\end{tabular}}
\end{center}

\loesung
\schritt{Schritt 1 --- Gewichtete Summen}
\[ A: \ 35\cdot1{,}00 + 33\cdot0{,}00 + 32\cdot0{,}40 = 35{,}0 + 0 + 12{,}8 = 47{,}8 \]
\[ B: \ 35\cdot0{,}90 + 33\cdot1{,}00 + 32\cdot0{,}00 = 31{,}5 + 33{,}0 + 0 = 64{,}5 \]
\[ C: \ 35\cdot0{,}00 + 33\cdot0{,}95 + 32\cdot1{,}00 = 0 + 31{,}35 + 32{,}0 = 63{,}4 \]

\schritt{Schritt 2 --- Das Ergebnis}
\begin{center}
{\small
\renewcommand{\arraystretch}{1.48}
\begin{tabular}{@{}P{2.4cm}P{2.2cm}P{2.6cm}@{}}
\toprule
& \textbf{Summe} & \textbf{Rang}\\
\midrule
$B$ & $64{,}5$ & $\mathbf{1}$\\
$C$ & $63{,}4$ & $2$\\
$A$ & $47{,}8$ & $3$\\
\bottomrule
\end{tabular}}
\end{center}

\begin{center}
\begin{minipage}{0.9\textwidth}
\textbf{Kein Zirkel --- aber $A$ landet auf dem letzten Platz, obwohl es die meisten ersten Plätze hatte.}
\end{minipage}
\end{center}

\schritt{Schritt 3 --- Woran das liegt}
\begin{center}
\begin{minipage}{0.88\textwidth}
\kursiv{\textbf{$A$ ist für Arzt II unerträglich ($0{,}00$), während $B$ für niemanden unerträglich ist.} Die Summe berücksichtigt die \kursiv{Stärke} der Vorlieben, die bloße Rangfolge nicht. Das ist der Gewinn --- und zugleich die Gefahr: Wer weiß, dass Stärken zählen, gibt extreme Werte an. Ein Verfahren, das dagegen unempfindlich ist, gibt es nach Band 21, Kapitel~13 nicht.}
\end{minipage}
\end{center}
\end{aufgabe}

\begin{geschichte}
\boxtitelg{Ein Marquis und ein Ökonom}
\Person{Nicolas de Condorcet} beschrieb den Zirkel 1785 in einer Schrift über die Anwendung der Wahrscheinlichkeitsrechnung auf Mehrheitsentscheidungen.

\begin{center}
\begin{minipage}{0.9\textwidth}
\kursiv{Er entwarf zugleich das nach ihm benannte Verfahren --- gewinnen soll, wer jeden Paarvergleich gewinnt --- und bewies dessen bekanntes Jury-Theorem: \textbf{Eine Mehrheit vieler mäßig zuverlässiger Urteile ist zuverlässiger als jedes einzelne.} Condorcet starb 1794 im Gefängnis der Revolution, deren Verfassung er mitentworfen hatte.}
\end{minipage}
\end{center}

\schritt{Arrow}
\begin{center}
\begin{minipage}{0.9\textwidth}
\kursiv{\Person{Kenneth Arrow} bewies den Unmöglichkeitssatz 1950 in seiner Doktorarbeit; er war damals achtundzwanzig. \textbf{1972 erhielt er als bis heute jüngster Preisträger den Wirtschaftsnobelpreis.} Der Satz beendete die Suche nach dem \enquote{richtigen} Wahlverfahren und verwandelte sie in die Frage, welche Bedingung man aufzugeben bereit ist.}
\end{minipage}
\end{center}
\end{geschichte}

\begin{forschung}
\boxtitelp{Gemeinsam entscheiden heute}
\begin{center}
{\small
\renewcommand{\arraystretch}{1.45}
\begin{tabular}{@{}P{3.4cm}P{4.8cm}@{}}
\toprule
\textbf{Strukturierte Konsensverfahren} & Delphi-Runden und Nominalgruppen in der Leitlinienarbeit.\\
\textbf{Urteilsaggregation} & Nicht Vorlieben, sondern Tatsachenurteile zusammenführen --- mit eigenen Unmöglichkeitssätzen.\\
\textbf{Gremiengüte} & Wann viele Urteile besser sind als wenige gute.\\
\bottomrule
\end{tabular}}
\end{center}

\schritt{Zur dritten Zeile}
\begin{center}
\begin{minipage}{0.88\textwidth}
\kursiv{Condorcets Jury-Theorem verlangt, dass die Urteile \kursiv{unabhängig} sind --- genau die Bedingung aus Band 22, Kapitel~9, unter der das Mitteln hilft. \textbf{In einer Konferenz, in der alle dieselbe Leitlinie gelesen haben und der Chefarzt zuerst spricht, ist sie verletzt.} Deshalb lassen gute Verfahren zunächst jeden schriftlich und verdeckt urteilen.}
\end{minipage}
\end{center}
\end{forschung}

\begin{merksatz}
\boxtitel{Was Sie jetzt können}
\begin{center}
\renewcommand{\arraystretch}{1.5}
\begin{tabular}{@{}P{3.0cm}P{5.6cm}@{}}
\toprule
\textbf{Condorcet-Zirkel} & Mehrheiten ohne Rangfolge\\
\textbf{Satz 13.1} & vier Bedingungen erzwingen den Diktator\\
\textbf{Medianwähler} & auf einer Linie gibt es einen Sieger\\
\textbf{Nutzen addieren} & Stärken zählen, Ehrlichkeit wird zur Frage\\
\textbf{Tagesordnung} & bestimmt das Ergebnis\\
\bottomrule
\end{tabular}
\end{center}

\begin{center}
\begin{minipage}{0.88\textwidth}
\kursiv{Bisher entschieden mehrere über dieselbe Sache. Das nächste Kapitel entscheidet über viele Menschen zugleich --- und dort wird die Summe zum Problem.}
\end{minipage}
\end{center}
\end{merksatz}
